\documentclass[conference]{IEEEtran}

\usepackage[utf8]{inputenc}
\usepackage[T1]{fontenc}
\usepackage[T1]{fontenc}
\usepackage{amsmath}
\usepackage{booktabs}
\usepackage{hyperref}
\usepackage{graphicx}
\usepackage{array}
\usepackage{microtype}

\hypersetup{
  colorlinks=true,
  linkcolor=blue,
  citecolor=blue,
  urlcolor=blue
}

% -------------------------------------------------------
\begin{document}

\title{Clasificación Automática de Preguntas sobre OEE\\
mediante Fine-Tuning de BERT en Español}

\author{
  \IEEEauthorblockN{E. Manosalvas}
  \IEEEauthorblockA{
    Maestría en Inteligencia Artificial\\
    Universidad San Francisco de Quito\\
    Quito, Ecuador\\
    \textit{emanosalvasb@estud.usfq.edu.ec}
  }
}

\maketitle

% -------------------------------------------------------
\begin{abstract}
Este artículo presenta el desarrollo y evaluación de un sistema de clasificación
automática de preguntas técnicas sobre el indicador de Efectividad Global del Equipo (OEE).
Se empleó el modelo BERT preentrenado en español
(\texttt{dccuchile/bert-base-spanish-wwm-cased}) como encoder contextual congelado,
combinado con una cabeza de clasificación lineal entrenada mediante fine-tuning
sobre un dataset balanceado de 524 preguntas en cuatro categorías:
\textit{calidad}, \textit{velocidad}, \textit{rendimiento} y \textit{eficiencia OEE}.
El modelo alcanzó una exactitud del 90.57\,\% en prueba y 92.31\,\% en validación,
con un F1-score macro de 0.921, demostrando la viabilidad de aplicar PLMs
al dominio industrial en idioma español.
\end{abstract}

\begin{IEEEkeywords}
BERT, clasificación de texto, OEE, procesamiento de lenguaje natural,
PyTorch Lightning, transfer learning, manufactura.
\end{IEEEkeywords}

% -------------------------------------------------------
\section{Introducción}

La Efectividad Global del Equipo (OEE) integra tres dimensiones clave del
desempeño productivo: disponibilidad, rendimiento de velocidad y calidad.
Diagnosticar cada dimensión a través de preguntas técnicas especializadas
es fundamental para sistemas de soporte de decisiones en manufactura.

Con la consolidación de los modelos de lenguaje preentrenados (PLMs) —en
particular BERT \cite{devlin2019bert}— se abrió la posibilidad de adaptar
representaciones contextuales a tareas de clasificación en dominios
especializados con datos reducidos. Sin embargo, la mayoría de los estudios
se desarrollan en inglés. Para el español, BETO \cite{canete2020spanish}
representa el estado del arte, pero su aplicación al dominio industrial del
OEE aún es limitada.

Este trabajo contribuye con: (1) un dataset balanceado de 524 preguntas
OEE en español; (2) un clasificador BERT+cabeza lineal entrenado con
PyTorch Lightning; y (3) una evaluación cuantitativa detallada por clase.

% -------------------------------------------------------
\section{Metodología}

\subsection{Dataset}

El ddataset consta de 524 preguntas técnicas en español distribuidas
de forma balanceada en cuatro clases de etiqueta única (\textit{single-label}),
con 131 ejemplos por clase (Tabla~\ref{tab:dataset}).
El dataset fue generado sintéticamente para cubrir el espacio semántico de cada
categoría. Una versión preliminar con esquema multi-label presentaba fuerte
solapamiento entre etiquetas, induciendo sesgo hacia la clase \textit{velocidad};
la transición a single-label balanceado corrigió este problema.
La partición se realizó en entrenamiento (80\,\%), validación (10\,\%) y
prueba (10\,\%) con semilla fija (\texttt{RANDOM\_SEED=42}).

\begin{table}[h]
\centering
\caption{Distribución del dataset por clase y partición}
\label{tab:dataset}
\begin{tabular}{lcccc}
\toprule
\textbf{Clase} & \textbf{Total} & \textbf{Train} & \textbf{Val} & \textbf{Test} \\
\midrule
calidad        & 131 & 105 & 13 & 13 \\
velocidad      & 131 & 105 & 13 & 13 \\
rendimiento    & 131 & 105 & 13 & 13 \\
eficiencia OEE & 131 & 103 & 13 & 15 \\
\midrule
\textbf{Total} & \textbf{524} & \textbf{418} & \textbf{52} & \textbf{54} \\
\bottomrule
\end{tabular}
\end{table}

\subsection{Arquitectura del Modelo}

El modelo sigue la estrategia de \textit{feature extraction}: el encoder
BERT-base en español (${\sim}109$\,M parámetros) se mantiene congelado en
modo \texttt{eval}, actuando exclusivamente como extractor de representaciones
contextuales. Sobre el vector del token \texttt{[CLS]} se añade una
cabeza lineal \texttt{Sequential} ($768\!\to\!4$, Dropout 0.1) con
3\,072 parámetros entrenables y \texttt{CrossEntropyLoss} como función
de pérdida. Esta decisión responde al tamaño moderado del dataset, donde
el fine-tuning completo del encoder sería propenso al sobreajuste.

\subsection{Entrenamiento}

El entrenamiento fue orquestado con PyTorch Lightning. Los principales
hiperparámetros se resumen en la Tabla~\ref{tab:hyper}. Se empleó
Early Stopping con paciencia de 5 épocas monitoreando
\texttt{val\_loss}, lo que detuvo el entrenamiento en la época 14
(checkpoint óptimo en época 13). El entrenamiento se ejecutó en GPU NVIDIA
con soporte CUDA.

\begin{table}[h]
\centering
\caption{Hiperparámetros de entrenamiento}
\label{tab:hyper}
\begin{tabular}{ll}
\toprule
\textbf{Hiperparámetro} & \textbf{Valor} \\
\midrule
Encoder            & bert-base-spanish-wwm-cased \\
Optimizador        & AdamW, lr $= 2\times10^{-5}$ \\
Scheduler          & Linear warmup (20\,\% pasos) \\
Batch size         & 32 \\
Max tokens         & 512 \\
Max épocas         & 50 (Early Stop: 5) \\
Pérdida            & CrossEntropyLoss \\
Dropout            & 0.1 \\
\bottomrule
\end{tabular}
\end{table}

% -------------------------------------------------------
\section{Resultados y Discusión}

\subsection{Métricas Globales}

El modelo alcanzó una exactitud de \textbf{90.57\,\%} en prueba y
\textbf{92.31\,\%} en validación, convergiendo en apenas 14 épocas.
Estos resultados son notables considerando que solo 3\,072 parámetros
fueron actualizados, lo que evidencia la potencia de las representaciones
contextuales de BERT para capturar diferencias semánticas en dominio
técnico especializado (Tabla~\ref{tab:global}).

\begin{table}[h]
\centering
\caption{Métricas globales en validación y prueba}
\label{tab:global}
\begin{tabular}{lcc}
\toprule
\textbf{Métrica} & \textbf{Validación} & \textbf{Prueba} \\
\midrule
Exactitud        & 92.31\,\% & 90.57\,\% \\
Pérdida (Loss)   & 0.4613    & 0.3656    \\
F1-score macro   & 0.9209    & ---       \\
F1-score ponderado & 0.9231  & ---       \\
\bottomrule
\end{tabular}
\end{table}

\subsection{Reporte por Clase}

La Tabla~\ref{tab:report} muestra el reporte de clasificación en validación
usando argmax. La clase \textit{calidad} obtiene precisión perfecta (1.000)
pero recall de 0.857, indicando que algunos ejemplos son asignados a otras
categorías. Por el contrario, \textit{eficiencia OEE} logra recall perfecto
(1.000) con precisión 0.857, sugiriendo tendencia a clasificar ejemplos limítrofes
en esta clase. La clase \textit{rendimiento} alcanza el F1 más alto (0.9375),
reflejando que sus representaciones BERT son las más discriminativas. Estos
resultados son consistentes con la literatura en clasificación de dominio
específico en español con BERT \cite{canete2020spanish}, donde se reportan
F1-scores de $0.89$--$0.95$.

\begin{table}[h]
\centering
\caption{Reporte de clasificación por clase — Validación (argmax)}
\label{tab:report}
\begin{tabular}{lcccc}
\toprule
\textbf{Clase} & \textbf{Prec.} & \textbf{Rec.} & \textbf{F1} & \textbf{Sup.} \\
\midrule
calidad        & 1.0000 & 0.8571 & 0.9231 & 14 \\
velocidad      & 0.9000 & 0.9000 & 0.9000 & 10 \\
rendimiento    & 0.9375 & 0.9375 & 0.9375 & 16 \\
eficiencia OEE & 0.8571 & 1.0000 & 0.9231 & 12 \\
\midrule
macro avg      & 0.9237 & 0.9237 & 0.9209 & 52 \\
weighted avg   & 0.9286 & 0.9231 & 0.9231 & 52 \\
\bottomrule
\end{tabular}
\end{table}

% -------------------------------------------------------
\section{Conclusiones}

Se demostró que un clasificador BERT con encoder congelado y cabeza lineal
alcanza resultados competitivos (test acc. 90.57\,\%, F1 macro 0.921) para
clasificar preguntas técnicas sobre OEE en español, entrenando únicamente
3\,072 parámetros. La transición de multi-label a single-label balanceado
fue determinante para eliminar sesgos sistemáticos en la predicción.

Entre las limitaciones destacan el uso de datos sintéticos y el tamaño
moderado del dataset. Como trabajo futuro se propone: ampliar el dataset con
preguntas reales de operadores; explorar fine-tuning completo con LoRA
\cite{hu2022lora}; e integrar el clasificador en un sistema QA para
entornos de manufactura reales.

% -------------------------------------------------------
\begin{thebibliography}{1}

\bibitem{devlin2019bert}
J.~Devlin, M.-W. Chang, K.~Lee, and K.~Toutanova,
``BERT: Pre-training of deep bidirectional transformers for language understanding,''
in \textit{Proc. NAACL-HLT}, Minneapolis, MN, USA, 2019, pp. 4171--4186.

\bibitem{canete2020spanish}
J.~Cañete, G.~Chaperon, R.~Fuentes, J.-H.~Ho, H.~Kang, and J.~Pérez,
``Spanish pre-trained BERT model and evaluation data,''
in \textit{PML4DC Workshop at ICLR}, 2020.

\bibitem{vaswani2017attention}
A.~Vaswani \textit{et al.},
``Attention is all you need,''
in \textit{Advances in Neural Information Processing Systems}, vol.~30, 2017.

\bibitem{hu2022lora}
E.~Hu \textit{et al.},
``LoRA: Low-rank adaptation of large language models,''
in \textit{Proc. ICLR}, 2022.

\end{thebibliography}

\end{document}
